% commentaire.tex — f166v, vv. 337–379 (conversion initiale des travaux antérieurs)
\comm{338–9}{ἢ τυραννεῖν ἐπινοεῖ ἢ τὸν τύραννον συγκατάγειν.}{Citation quasi-textuelle d'une ancienne loi, peut-être de Dracon, rapportée dans la \italique{Constitution d'Athènes} (XVI, 10) :
\citComm{ἦσαν δὲ καὶ τοῖς Ἀθηναίοις οἱ περὶ τῶν τυράννων νόμοι πρᾷοι κατ᾽ ἐκείνους τοὺς καιρούς, οἵ τ᾽ ἄλλοι καὶ δὴ καὶ ὁ μάλιστα καθήκων πρὸς τὴν τῆς τυραννίδος <κατάστασιν>. νόμος γὰρ αὐτοῖς ἦν ὅδε. `θέσμια τάδε Ἀθηναίων καὶ πάτρια· ἐάν τινες τυραννεῖν ἐπανιστῶνται ἐπὶ τυραννίδι, ἢ συγκαθιστῇ τὴν τυραννίδα, ἄτιμον εἶναι καὶ αὐτὸν καὶ γένος.'\\
\francais{Les lois des Athéniens sur la tyrannie étaient alors [à l'époque de Pisitrate] peu sévères, celle surtout qui visait les entreprises des tyrans et dont voici le texte : « Les lois athéniennes établies par nos pères portent que : quiconque aspire à la tyrannie ou forme un complot pour l'établir, sera frappé d'atimie dans sa personne et dans sa race. »}}Étant donnée l'époque troublée de la représentation des \italique{Thesmophories}, cette loi devait être sur toutes les bouches ; pour preuve, la même année que la représentation des \italique{Thesmophories}, en 410/09\nc{C. Carey, Ifigeneia G., Brenda G.-W., \italique{Use and abuse of law in the Athenian courts}, p. 295.}, Démophante proposa un décret inspiré de cette dernière garantissant l'impunité à l'assassin de toute personne complice de crime de lèse-démocratie (cette loi fut gravée sur une stèle fameuse à Athènes qu'Andocide cite dans son \italique{Discours sur les mystères}, 97) :\citComm{\grec{καὶ ἐάν τις τυραννεῖν ἐπαναστῇ ἢ τὸν τύραννον συγκαταστήσῃ}.}}
\comm{346}{ἢ καὶ δέχεται προδιδοῦσ᾽ ἑταίρᾳ τὸν φίλον.}{Pour le texte (ἑταίρᾳ) et le sens du vers (la vieille entremetteuse), voir le commentaire critique.}
\comm{347–8}{κεἴ τις κάπηλος ἢ καπηλὶς τοῦ χοὸς ἢ τῶν κοτυλῶν τὸ νόμισμα διαλυμαίνεται.}{Les débitantes de mauvais aloi sont brocardées de la même façon par Aristophane dans le \italique{Ploutos}, v. 435-6 :
\citComm{Βλεψίδημος — ἆρ᾽ ἐστὶν ἡ καπηλὶς ἡ ᾽κ τῶν γειτόνων,\\
ἣ ταῖς κοτύλαις ἀεί με διαλυμαίνεται;}}
\comm{356–67}{}{Pour le texte et l'interprétation (\nieme{10} note de Willems), voir le commentaire critique.}
\comm{365–6}{ἢ Μήδους ἐπάγουσι τῆς χώρας / οὕνεκ᾽ ἐπὶ βλάβῃ.}{Pour le texte, voir le commentaire critique, qui donne le début de l'extrait de l'article de Taillardat \& Blanc ; en voici la suite : Il y a cependant autre chose dans ce texte : un jeu de mots sur τῆς χώρας, plaisanterie dont personne ne s'est encore avisé. En effet, le parallélisme éclatant, et sûrement dû à Aristophane lui-même, de κερδῶν οὕνεκ᾽ ἐπὶ βλάβῃ (v. 360) et de τῆς χώρας οὕνεκ᾽ ἐπὶ βλάβῃ (v. 366) incite à chercher dans τῆς χώρας, autant que le nom de l'Attique, une variation de la notion exprimée par κερδῶν ; car le chœur peut faire allusion, ici comme là, à l'intérêt personnel et à l'avidité des gens qu'il maudit. Or le mot χώρα désigne, chez plusieurs auteurs, “la place” qu'on occupe dans la vie, par exemple “le rang”, qu'on a dans la hiérarchie sociale ou militaire\nc{Xén., \italique{Anab.}, 5, 6, 13 (ἐν ἀνδραπόδων χώρᾳ ἐσόμεθα), \italique{Cyr.}, 2, 1, 18 (ἐν μισθοφόρου χώρᾳ εἶναι), Polybe, 1, 43, 1 (τῶν ἡγεμόνων τινὲς τῶν τὰς μεγίστας χώρας ἐχόντων ἐν τοῖς μισθοφόροις), Pausanias, 10, 1, 8 (ὁ δὲ χώραν ἐν τοῖς ἄρχουσιν ἔχων τὴν μεγίστην μάντις ἦν Τελλίας ὁ Ἠλεῖος). Avec le sens de « place dans la société », χώρα paraît déjà exister chez Théognis, v. 822.}. Dans les \italique{Thesmophories}, v. 366, c'est de rang dans la hiérarchie politique et sociale qu'il s'agit ; si bien que τῆς χώρας οὕνεκα, signifie à la fois “pour l'amour de notre pays (disent-elles)” et “pour se pousser aux (bonnes) places”\nc{On justifiera ci-dessous la traduction « aux \italique{bonnes} places ».}, littéralement “en vue de leur place personnelle”\nc{Tel est le sens de l'article, cf. Kühner-Gerth, \italique{Satzlehre}, I, p. 593, n°2. C'est l'emploi intentionnel du singulier τῆς χώρας (Kühner-Gerth, I, p. 14, n°4), au lieu du pluriel τῶν χωρῶν (plus fréquent dans ce cas, cf. Kühner-Gerth, I, p. 19, Anm. 3), qui permet le jeu de mots. Pour οὕνεκα, \cf \italique{Lysistrata}, v. 491 : τοῦδ᾽ οὕνεκα « en vue de cela (= en vue des magistratures et des vols qu'elles permettent) ».}. 
Plus que le cas général des σπουδαρχίδαι, “les sires de Brigue-charges”\nc{\italique{Ach.}, 595 ; pour le fait, \cf \italique{Equ.}, 157-159, 164-165, et \italique{Aves}, 798-800.}, Aristophane vise ici une circonstance particulière. Il pense, comme Willems l'a dit, aux intrigues ambitieuses d'Alcibiade\nc{Envisageant de provoquer une révolution à Athènes, il préparait alors son retour dans sa patrie (Thuc., 8, 47, 1-2).} et de Pisandre qui, pendant l'hiver 412-411, prétendaient arracher la Perse à son alliance spartiate et cherchaient à obtenir, pour Athènes, l'appui et l'or du Grand Roi\nc{Pour le rôle de l'or perse à ce moment, voir Thucydide, 8, 29, 1-2 ; 8, 45, 2-6; 8, 48, 2-3 ; 8, 53, 2.}, pour eux-mêmes, le pouvoir à Athènes. Et que fait Pisandre, à la fin de 412, quand le peuple le délègue auprès d'Alcibiade et de Tissapherne ? Avant de quitter l'Attique pour l'Asie Mineure, “il rend visite aux sociétés secrètes déjà établies à Athènes pour s'occuper des procès et des magistratures”, ξυνωμοσίας, αἵπερ ἐτύγχανον πρότερον ἐν τῇ πόλει οὖσαι ἐπὶ δίκαις καὶ ἀρχαῖς, ἁπάσας ἐπελθὼν... (Thuc. VIII, 54, 4). C'est dire que les comploteurs avaient ouvert la chasse aux places. Les oligarques se consultent ἐπὶ δίκαις καὶ ἀρχαῖς et tout le monde le sait à Athènes : deux mois avant les \italique{Thesmophories}, Aristophane a déjà dénoncé “ces gens qui s'agglomèrent et font bourre pour se pousser aux charges” τούς γε συνισταμένους τούτους καὶ τοὺς πιλοῦντας ἑαυτοὺς | ἐπὶ ταῖς ἀρχαῖσι (\italique{Lysistrata}, v. 577 sq.). Ces ambitieux sont concussionnaires et vénaux : “c'est pour être en mesure de voler que Pisandre et ceux qui ne pensent qu'aux charges créent à chaque instant grabuge et désordre. Qu'ils fassent donc, pour cela, tout ce qu'ils voudront ; mais cet argent-ci\nc{Le trésor déposé à l'Acropole dont les Athéniennes barricadées interdisent l'entrée.}, il n'y a plus de danger qu'ils l'emportent”, 
\citComm{ἵνα γὰρ Πείσανδρος ἔχοι κλέπτειν χοἰ ταῖς ἀρχαῖς ἐπέχοντες,\\
ἀεί τινα κορκορυγὴν ἐκύκων. οἱ δ᾽ οὖν τοῦδ᾽ οὕνεκα δρώντων\\
ὅ τι βούλονται· τὸ γὰρ ἀργύριον τοῦτ᾽ οὐκέτι μὴ καθέλωσιν.\\
(\italique{Lysistrata}, v. 490-492).}Ἐπὶ (δίκαις καὶ) ἀρχαῖς, dit Thucydide ; ἐπὶ ταῖς ἀρχαῖσι, οἱ ταῖς ἀρχαῖς ἐπέχοντες, dit Aristophane. Peut-on rêver plus belles gloses au τῆς χώρας οὕνεκα des \italique{Thesmophories} ? Aux yeux d'Aristophane, tremper dans la grenouillère des combinaisons louches de l'hiver de 412-411, c'est donc prétendre à faire d'une pierre deux coups : devenir un personnage important à Athènes et s'enrichir en profitant de tous les avantages, licites ou non\nc{Voir J. Taillardat, \italique{Les images d'Aristophane}, pp. 408-423. Pour les avantages licites, \cf \italique{Acharniens}, v. 65 sqq. (indemnité de mission officielle et vie facile des ambassadeurs athéniens auprès du Grand Roi).}, attachés aux ἀρχαί ; c'est tout cela que résume τῆς χώρας οὕνεκ᾽ ἐπὶ βλάβῃ, car de telles ambitions ne réussissent qu'au mépris du bien public.}
\comm{376–7}{ἐκκλησίαν ποιεῖν ἕωθεν τῇ μέσῃ / τῶν Θεσμοφορίων, ᾗ μάλισθ᾽ ἡμῖν σχολή.}{Remarquons au passage que les assemblées se tenaient à l'aube\nc{Coutume confirmée par d'autres passages d'Aristophane : \italique{Acharniens}, 19-20 ; \italique{Assemblée des femmes}, 20-21, 84-85, 739-741.}. Le fait que le \italique{jour du milieu} soit celui du \italique{jeûne} (νηστεία) explique que les femmes aient tout le loisir pour tenir assemblée.}
\comm{373–9}{ἄκουε πᾶς ... τίς ἀγορεύειν βούλεται;}{L'assemblée à huis-clos féminin se donne les formes pour être en droit de juger Euripide : tout ce passage est un pastiche d'un \grec{προβούλευμα}, une \italique{résolution du Conseil}. En voici un exemple : le début du décret proposé par Démophante évoqué plus haut (Andocide, \italique{Discours sur les mystères}, 96) :
\citComm{ΝΟΜΟΣ\\
Ἔδοξε τῇ βουλῇ καὶ τῷ δήμῳ. Αἰαντὶς ἐπρυτάνευε, Κλειγένης ἐγραμμάτευε, Βοηθὸς ἐπεστάτει. Τάδε Δημόφαντος συνέγραψεν. Ἄρχει χρόνος τοῦδε τοῦ ψηφίσματος ἡ βουλὴ οἱ πεντακόσιοι λαχόντες τῷ κυάμῳ, οἷς Κλειγένης πρῶτος ἐγραμμάτευεν.\nc{Autres exemples : Thucydide, IV, 118 ; \italique{Inscriptiones Graecae} I\expo{3} 5, 9, 10, 12, 17 \&c. \&c.}}}
